\documentclass[12pt]{article}
\makeatletter
\def\input@path{{../../}}
\makeatother
\usepackage{sbc-template}
\makeatletter
\renewcommand{\inst}[1]{}
\makeatother
\usepackage{graphicx,url}
\usepackage[utf8]{inputenc}
\usepackage[T1]{fontenc}
\usepackage[brazil]{babel}
\usepackage{longtable,booktabs,array}
\providecommand{\tightlist}{%%
  \setlength{\itemsep}{0pt}\setlength{\parskip}{0pt}}
\sloppy
\title{Por que Reservoir Computing funciona: memoria, nao linearidade e capacidade de representacao}
\author{}
\address{}
\begin{document}

\maketitle

\begin{resumo}
Depois do primeiro experimento reproduzivel, a pergunta natural deixa de ser "o RC ganhou?" e passa a ser "o que exatamente o RC esta tentando fazer?". Este artigo responde a essa pergunta em tres eixos: memoria, nao linearidade e capacidade de representacao. Em vez de tratar o reservatorio como uma caixa-preta, mostramos como hiperparametros como \texttt{n\_reservoir}, \texttt{spectral\_radius}, \texttt{leak\_rate} e \texttt{washout} alteram o comportamento do modelo. O objetivo nao e vender RC como solucao universal, mas oferecer uma leitura mecanica do metodo. Ao final, o leitor entende por que um reservatorio pode capturar dependencias temporais e misturas nao lineares e tambem por que isso, sozinho, ainda nao garante superioridade sobre baselines fortes.
\end{resumo}

\section{1. Introducao}\label{1-introducao}

No artigo anterior, o RC nao venceu o Seasonal Naive no recorte inicial do Favorita. Esse resultado pode ser lido de duas formas. A leitura superficial diria: "o metodo nao funcionou". A leitura interessante diz: "agora temos uma oportunidade de entender o metodo".

Reservoir Computing merece ser estudado nao apenas porque pode performar bem, mas porque organiza de forma muito clara uma pergunta central em modelagem temporal: como transformar uma sequencia em uma representacao dinamica que preserve informacao util sobre o passado?

Para responder a isso, precisamos abrir a caixa do reservatorio.

\section{2. Memoria: o passado precisa continuar presente}\label{2-memoria-o-passado-precisa-continuar-presente}

Em forecasting, memoria significa que a previsao de hoje depende de sinais que nao estao apenas na observacao atual. No Favorita, isso aparece de modo direto:

\begin{itemize}
\tightlist
\item
  a demanda de um dia carrega padroes do dia da semana;
\item
  promocao atual pode herdar efeito da promocao anterior;
\item
  o comportamento perto de feriados costuma depender do contexto recente.
\end{itemize}

O reservatorio tenta condensar esse historico em um estado interno \texttt{x\_t}. Esse estado nao e um buffer literal dos ultimos dias. Ele e uma compressao dinamica do passado.

\subsection{2.1 O papel do leak rate}\label{21-o-papel-do-leak-rate}

No projeto, o RC usa \texttt{leak\_rate\ =\ 0.15}. Em termos intuitivos:

\begin{itemize}
\tightlist
\item
  leak rate baixo faz o estado mudar mais devagar e guardar mais passado;
\item
  leak rate alto faz o estado reagir mais rapidamente ao presente.
\end{itemize}

Isso cria um trade-off importante. Se a dinamica ficar lenta demais, o modelo pode suavizar em excesso. Se ficar rapida demais, pode esquecer o passado cedo demais.

\subsection{2.2 O papel do washout}\label{22-o-papel-do-washout}

O \texttt{washout\ =\ 14} usado no projeto lembra que memoria util nao e o mesmo que inicializacao arbitraria. Antes de treinar o readout, damos tempo para que o reservatorio entre em um regime ditado pelos dados e nao pelo estado inicial.

\section{3. Nao linearidade: por que uma serie de demanda nao cabe em uma regra unica}\label{3-nao-linearidade-por-que-uma-serie-de-demanda-nao-cabe-em-uma-regra-unica}

Muitos padroes do varejo nao sao lineares. Um mesmo aumento de promocao nao produz sempre o mesmo efeito. O impacto de um sabado pode depender da epoca do mes. A proximidade de um feriado pode amplificar um comportamento que, em outra semana, seria pequeno.

A nao linearidade do reservatorio, tipicamente implementada com \texttt{tanh}, ajuda a produzir uma resposta mais rica a combinacoes de sinais. Ela nao "resolve" toda a complexidade do problema, mas permite que o estado interno misture historico e contexto de forma menos rigida do que um modelo linear puro.

No projeto, essa nao linearidade atua sobre um vetor de entrada pequeno e interpretavel:

\begin{itemize}
\tightlist
\item
  demanda previa;
\item
  promocao;
\item
  codificacoes ciclicas de calendario;
\item
  indicador de fim de semana.
\end{itemize}

Essa escolha e pedagogica porque mostra que nao e preciso comecar com centenas de features para discutir representacao temporal.

\section{4. Capacidade de representacao: por que tamanho importa, mas nao sozinho}\label{4-capacidade-de-representacao-por-que-tamanho-importa-mas-nao-sozinho}

Um reservatorio maior tende a oferecer um espaco de estados mais rico. No entanto, mais unidades nao significam automaticamente melhor desempenho. O tamanho do reservatorio controla quanta diversidade de resposta dinamica o modelo pode sustentar, mas essa diversidade so ajuda quando o readout consegue explorá-la e quando ela e coerente com o problema.

No projeto, o RC atual usa \texttt{n\_reservoir\ =\ 80}. Esse valor foi escolhido como um compromisso entre:

\begin{itemize}
\tightlist
\item
  simplicidade didatica;
\item
  custo computacional moderado;
\item
  capacidade suficiente para um primeiro experimento.
\end{itemize}

Essa decisao reforca um ponto importante para o leitor: hiperparametro em RC nao e decoracao. Ele expressa uma hipotese sobre a dinamica necessaria para o problema.

\section{5. Como ler os hiperparametros do RC}\label{5-como-ler-os-hiperparametros-do-rc}

Uma boa forma de ensinar RC e mostrar o que cada hiperparametro faz quando esta "baixo demais", "equilibrado" ou "alto demais". A tabela abaixo resume essa leitura.

\begin{longtable}[]{@{}llll@{}}
\toprule\noalign{}
Hiperparametro & Muito baixo & Faixa equilibrada & Muito alto \\
\midrule\noalign{}
\endhead
\bottomrule\noalign{}
\endlastfoot
\texttt{n\_reservoir} & pouca diversidade de estados & representacao rica e controlavel & custo maior e risco de redundancia \\
\texttt{spectral\_radius} & dinamica curta e muito contrativa & memoria estavel e resposta temporal rica & risco de instabilidade e excesso de persistencia \\
\texttt{leak\_rate} & memoria longa, resposta lenta & equilibrio entre passado e presente & reatividade alta e pouca memoria \\
\texttt{washout} & estado inicial contamina treino & treino sobre dinamica estabilizada & descarte desnecessario de dados \\
\texttt{ridge\_alpha} & risco de sobreajuste no readout & regularizacao moderada & leitura excessivamente rigida \\
\end{longtable}

Essa tabela substitui bem, para fins didaticos, a ideia vaga de "fazer tuning". Em vez de pensar tuning como busca cega, o leitor passa a pensar tuning como ajuste de propriedades dinamicas.

\section{6. O que o resultado atual do RC sugere}\label{6-o-que-o-resultado-atual-do-rc-sugere}

No benchmark atual do projeto, o RC obteve:

\begin{itemize}
\tightlist
\item
  \texttt{MAE\ =\ 324.294}
\item
  \texttt{RMSE\ =\ 423.417}
\item
  \texttt{WAPE\ =\ 0.1498}
\item
  \texttt{sMAPE\ =\ 0.1653}
\end{itemize}

Esses numeros colocam o RC acima da LSTM simples do projeto, mas abaixo de ETS, XGBoost, Prophet e Seasonal Naive.

Essa posicao intermediaria sugere uma leitura produtiva:

\begin{enumerate}
\def\labelenumi{\arabic{enumi}.}
\tightlist
\item
  o reservatorio captura algo relevante, pois nao fracassa completamente;
\item
  o pipeline ainda nao extrai o maximo valor da representacao temporal;
\item
  modelos classicos fortes continuam sendo referencias serias para previsao de demanda.
\end{enumerate}

Em outras palavras, o RC ensina mesmo quando nao vence.

\section{7. O que RC faz bem e o que ainda nao faz bem neste problema}\label{7-o-que-rc-faz-bem-e-o-que-ainda-nao-faz-bem-neste-problema}

\subsection{7.1 O que RC faz bem}\label{71-o-que-rc-faz-bem}

\begin{itemize}
\tightlist
\item
  integra historico recente e contexto exogeno em um unico estado;
\item
  produz uma representacao dinamica reutilizavel por um readout simples;
\item
  oferece uma linguagem clara para falar de memoria temporal.
\end{itemize}

\subsection{7.2 O que RC ainda nao faz bem aqui}\label{72-o-que-rc-ainda-nao-faz-bem-aqui}

\begin{itemize}
\tightlist
\item
  nao supera uma sazonalidade semanal forte quando o baseline ja a captura muito bem;
\item
  ainda opera em um recorte pequeno e com feature set enxuto;
\item
  depende de tuning cuidadoso para explorar melhor a estrutura da serie.
\end{itemize}

Essa avaliacao e especialmente importante porque prepara a ponte para QRC. Se o leitor nao entender RC como paradigma, QRC parecerá apenas uma troca de implementacao. Se entender, percebera continuidade entre os dois.

\section{8. Uma leitura didatica da capacidade de representacao}\label{8-uma-leitura-didatica-da-capacidade-de-representacao}

Uma forma util de resumir o artigo e a seguinte:

\begin{itemize}
\tightlist
\item
  memoria responde a pergunta "quanto do passado o modelo ainda carrega?";
\item
  nao linearidade responde a pergunta "como sinais se combinam de forma nao aditiva?";
\item
  capacidade de representacao responde a pergunta "quantas respostas dinamicas diferentes o sistema consegue sustentar?".
\end{itemize}

Essas tres dimensoes sao suficientes para interpretar boa parte do comportamento de um reservatorio.

\section{9. Conclusao}\label{9-conclusao}

Reservoir Computing funciona quando seu estado interno consegue reter passado, reagir de forma nao linear e fornecer uma representacao util para o readout. Isso nao garante vitoria em benchmark, mas fornece um quadro conceitual extremamente forte para pensar series temporais.

Esse quadro agora nos permite dar um passo essencial: comparar RC com outras familias de modelos de forma justa, rigorosa e reproduzivel.

\section{Entregaveis associados no repositorio}\label{entregaveis-associados-no-repositorio}

\begin{itemize}
\tightlist
\item
  implementacao do RC: \texttt{code/rc/model.py}
\item
  resultados do RC: \texttt{code/rc/RESULTS.md}
\item
  benchmark consolidado: \texttt{code/RESULTS.md}
\end{itemize}

\section{Referencias}\label{referencias}

\begin{enumerate}
\def\labelenumi{\arabic{enumi}.}
\tightlist
\item
  Jaeger, H. The "echo state" approach to analysing and training recurrent neural networks.
\item
  Lukosevicius, M., Jaeger, H. Reservoir computing approaches to recurrent neural network training.
\item
  Maass, W., Natschlager, T., Markram, H. Real-time computing without stable states.
\end{enumerate}

\section*{Resultados Computacionais Associados}

Os resultados computacionais associados a este artigo estao organizados na subpasta \texttt{computational\_results\_20260402\_222902/}, localizada dentro desta pasta \LaTeX. Esse diretorio contem o arquivo \texttt{REPORT.md}, tabelas em CSV e imagens comparativas apropriadas ao escopo do artigo.

\end{document}
